\textbf{Phần 1. Trắc nghiệm lựa chọn}
\setcounterpageref{ex}{0}

\begin{ex} %[3D3V30]
Điền số thích hợp vào chỗ chấm: $5\text{ cm } 8\text{ mm} = \dots\text{ mm}$.
\choice
{$508\text{ mm}$}
{$580\text{ mm}$}
{$50\text{ mm}$}
{\True $58\text{ mm}$}
\loigiai{
Vì $1\text{ cm} = 10\text{ mm}$ nên $5\text{ cm } 8\text{ mm} = 50\text{ mm} + 8\text{ mm} = 58\text{ mm}$.
\textbf{Đáp án của câu hỏi là D}
}
\end{ex}

\begin{ex} %[3D3V30]
Một tập giấy dày $24\text{ mm}$ gồm 120 trang giấy giống nhau. Hỏi một tập giấy tương tự dày $48\text{ mm}$ sẽ có bao nhiêu trang giấy?
\choice
{$180$ trang}
{\True $240$ trang}
{$300$ trang}
{$200$ trang}
\loigiai{
Tập giấy dày $48\text{ mm}$ gấp số lần tập $24\text{ mm}$ là: $48 : 24 = 2$ (lần).\
Số trang giấy là: $120 \times 2 = 240$ (trang).
\textbf{Đáp án của câu hỏi là B}
}
\end{ex}

\begin{ex} %[3D3V30]
Một con ốc sên mỗi phút bò được $15\text{ mm}$. Hỏi sau 6 phút con ốc sên bò được đoạn đường dài bao nhiêu xăng-ti-mét?
\choice
{$90\text{ cm}$}
{$6\text{ cm}$}
{\True $9\text{ cm}$}
{$15\text{ cm}$}
\loigiai{
Quãng đường ốc sên bò được: $15 \times 6 = 90\text{ (mm)}$. Đổi $90\text{ mm} = 9\text{ cm}$.
\textbf{Đáp án của câu hỏi là C}
}
\end{ex}

\begin{ex} %[3D3V30]
Một sợi dây đồng dài $1\text{ m}$. Người ta cắt đi một đoạn dài $35\text{ cm } 6\text{ mm}$. Đoạn dây còn lại dài bao nhiêu mi-li-mét?
\choice
{\True $644\text{ mm}$}
{$654\text{ mm}$}
{$640\text{ mm}$}
{$744\text{ mm}$}
\loigiai{
Đổi $1\text{ m} = 1000\text{ mm}$; $35\text{ cm } 6\text{ mm} = 356\text{ mm}$.\
Đoạn dây còn lại là: $1000 - 356 = 644\text{ (mm)}$.
\textbf{Đáp án của câu hỏi là A}
}
\end{ex}

\begin{ex} %[3D3V30]
Một chiếc đinh dài $4\text{ cm } 5\text{ mm}$ được đóng ngập vào thanh gỗ một đoạn $32\text{ mm}$. Phần đinh còn nhô ra ngoài thanh gỗ dài bao nhiêu mi-li-mét?
\choice
{$15\text{ mm}$}
{$12\text{ mm}$}
{$18\text{ mm}$}
{\True $13\text{ mm}$}
\loigiai{
Đổi $4\text{ cm } 5\text{ mm} = 45\text{ mm}$. Phần đinh nhô ra: $45 - 32 = 13\text{ (mm)}$.
\textbf{Đáp án của câu hỏi là D}
}
\end{ex}

\begin{ex} %[3D3V30]
Một hình tam giác có độ dài ba cạnh lần lượt là $25\text{ mm}, 35\text{ mm}$ và $4\text{ cm}$. Chu vi hình tam giác đó theo xăng-ti-mét là
\choice
{$100\text{ cm}$}
{\True $10\text{ cm}$}
{$8\text{ cm}$}
{$12\text{ cm}$}
\loigiai{
Đổi $4\text{ cm} = 40\text{ mm}$. Chu vi là: $25 + 35 + 40 = 100\text{ mm} = 10\text{ cm}$.
\textbf{Đáp án của câu hỏi là B}
}
\end{ex}

\begin{ex} %[3D3V30]
Người ta xếp chồng 8 cuốn sách giống hệt nhau thì được một chồng sách cao $9\text{ cm } 6\text{ mm}$. Độ dày của mỗi cuốn sách là
\choice
{$14\text{ mm}$}
{$10\text{ mm}$}
{\True $12\text{ mm}$}
{$16\text{ mm}$}
\loigiai{
Đổi $9\text{ cm } 6\text{ mm} = 96\text{ mm}$. Độ dày mỗi cuốn sách: $96 : 8 = 12\text{ (mm)}$.
\textbf{Đáp án của câu hỏi là C}
}
\end{ex}

\begin{ex} %[3D3V30]
So sánh hai độ dài sau: $A = 7\text{ cm } 4\text{ mm}$ và $B = 740\text{ mm} : 10$. Khẳng định đúng là
\choice
{\True $A = B$}
{$A > B$}
{$A < B$}
{Không so sánh được}
\loigiai{
$A = 74\text{ mm}$, $B = 740 : 10 = 74\text{ mm}$. Vậy $A = B$.
\textbf{Đáp án của câu hỏi là A}
}
\end{ex}

\begin{ex} %[3D3V30]
Một cây bút chì ban đầu dài $18\text{ cm}$. Sau một thời gian sử dụng, bút chì bị gọt ngắn đi $65\text{ mm}$. Chiều dài hiện tại của cây bút chì là
\choice
{$12\text{ cm } 5\text{ mm}$}
{$11\text{ cm}$}
{$12\text{ cm}$}
{\True $11\text{ cm } 5\text{ mm}$}
\loigiai{
Đổi $18\text{ cm} = 180\text{ mm}$. Chiều dài còn lại: $180 - 65 = 115\text{ mm} = 11\text{ cm } 5\text{ mm}$.
\textbf{Đáp án của câu hỏi là D}
}
\end{ex}

\begin{ex} %[3D3V30]
Một con kiến càng dài $12\text{ mm}$, con ve sầu dài gấp 3 lần con kiến càng. Con ve sầu dài hơn con kiến càng bao nhiêu mi-li-mét?
\choice
{$36\text{ mm}$}
{\True $24\text{ mm}$}
{$20\text{ mm}$}
{$18\text{ mm}$}
\loigiai{
Ve sầu dài: $12 \times 3 = 36\text{ mm}$. Dài hơn kiến càng: $36 - 12 = 24\text{ (mm)}$.
\textbf{Đáp án của câu hỏi là B}
}
\end{ex}

\begin{ex} %[3D3V30]
Một hình vuông có chu vi là $20\text{ cm}$. Độ dài cạnh của hình vuông đó tính theo mi-li-mét là
\choice
{$5\text{ mm}$}
{$500\text{ mm}$}
{\True $50\text{ mm}$}
{$25\text{ mm}$}
\loigiai{
Cạnh hình vuông là: $20 : 4 = 5\text{ cm}$. Đổi $5\text{ cm} = 50\text{ mm}$.
\textbf{Đáp án của câu hỏi là C}
}
\end{ex}

\begin{ex} %[3D3V30]
Người ta nối 5 thanh nẹp gỗ bằng nhau, mỗi thanh dài $18\text{ cm } 4\text{ mm}$ thành một thanh dài (mối nối không đáng kể). Chiều dài thanh nẹp thu được là
\choice
{\True $92\text{ cm}$}
{$90\text{ cm}$}
{$94\text{ cm}$}
{$920\text{ cm}$}
\loigiai{
Đổi $18\text{ cm } 4\text{ mm} = 184\text{ mm}$. Chiều dài thanh thu được: $184 \times 5 = 920\text{ mm} = 92\text{ cm}$.
\textbf{Đáp án của câu hỏi là A}
}
\end{ex}

\textbf{Phần 2. Trắc nghiệm đúng sai}
\setcounterpageref{ex}{0}

\begin{ex} %[3D3V30]
Xét tính đúng sai của các mối quan hệ giữa mi-li-mét và các đơn vị đo độ dài khác:
\choiceTF
{\True $1\text{ cm} = 10\text{ mm}$.}
{$1\text{ m} = 100\text{ mm}$.}
{\True $1\text{ dm} = 100\text{ mm}$.}
{$3\text{ cm } 5\text{ mm} = 305\text{ mm}$.}
\loigiai{
\begin{itemize}
    \item Ý a) Đúng: Mối quan hệ chuẩn.
    \item Ý b) Sai: Sai vì $1\text{ m} = 1000\text{ mm}$.
    \item Ý c) Đúng: Vì $1\text{ dm} = 10\text{ cm} = 100\text{ mm}$.
    \item Ý d) Sai: Sai vì $3\text{ cm } 5\text{ mm} = 35\text{ mm}$.
\end{itemize}
\begin{itemize}
    \item a) \textbf{Đúng}.
    \item b) \textbf{Sai}.
    \item c) \textbf{Đúng}.
    \item d) \textbf{Sai}.
\end{itemize}
}
\end{ex}

\begin{ex} %[3D3V30]
Một thanh kim loại dài $12\text{ cm } 8\text{ mm}$ được cưa thành 4 đoạn bằng nhau:
\choiceTF
{\True Thanh kim loại ban đầu dài $128\text{ mm}$.}
{\True Mỗi đoạn dài $32\text{ mm}$.}
{Người thợ phải cưa đúng 4 nhát cưa.}
{\True Hai đoạn cộng lại dài $6\text{ cm } 4\text{ mm}$.}
\loigiai{
\begin{itemize}
    \item Ý a) Đúng: $12\text{ cm } 8\text{ mm} = 128\text{ mm}$.
    \item Ý b) Đúng: $128 : 4 = 32\text{ mm}$.
    \item Ý c) Sai: Để chia thành 4 đoạn chỉ cần $4 - 1 = 3$ nhát cưa.
    \item Ý d) Đúng: $32 \times 2 = 64\text{ mm} = 6\text{ cm } 4\text{ mm}$.
\end{itemize}
\begin{itemize}
    \item a) \textbf{Đúng}.
    \item b) \textbf{Đúng}.
    \item c) \textbf{Sai}.
    \item d) \textbf{Đúng}.
\end{itemize}
}
\end{ex}

\textbf{Phần 3. Trả lời ngắn}
\setcounterpageref{ex}{0}

\begin{ex} %[3D3V30]
Đổi $7\text{ cm } 9\text{ mm}$ ra đơn vị mi-li-mét.
\shortans{79}
\loigiai{
\textbf{Đáp số:} $79$\\
$7\text{ cm } 9\text{ mm} = 70 + 9 = 79\text{ mm}$.
}
\end{ex}

\begin{ex} %[3D3V30]
Một cuốn sổ tay dày $18\text{ mm}$. Hỏi 5 cuốn sổ tay như thế xếp chồng lên nhau thì dày bao nhiêu mi-li-mét?
\shortans{90}
\loigiai{
\textbf{Đáp số:} $90$\\
$18 \times 5 = 90\text{ mm}$.
}
\end{ex}

\begin{ex} %[3D3V30]
Một đoạn dây kẽm dài $6\text{ cm } 5\text{ mm}$ được uốn thành một hình ngũ giác đều (5 cạnh bằng nhau). Hỏi mỗi cạnh dài bao nhiêu mi-li-mét?
\shortans{13}
\loigiai{
\textbf{Đáp số:} $13$\\
$65 : 5 = 13\text{ mm}$.
}
\end{ex}

\begin{ex} %[3D3V30]
Tính giá trị: $45\text{ mm} \times 2 - 18\text{ mm}$ theo đơn vị mi-li-mét.
\shortans{72}
\loigiai{
\textbf{Đáp số:} $72$\\
$90 - 18 = 72\text{ mm}$.
}
\end{ex}

\begin{ex} %[3D3V30]
Một tấm tôn dày $3\text{ mm}$. Hỏi chồng 25 tấm tôn như thế có độ dày là bao nhiêu mi-li-mét?
\shortans{75}
\loigiai{
\textbf{Đáp số:} $75$\\
$3 \times 25 = 75\text{ mm}$.
}
\end{ex}

\begin{ex} %[3D3V30]
Chiều dài của một chiếc kẹp giấy là $28\text{ mm}$. Nối 3 chiếc kẹp giấy móc vào nhau dài $76\text{ mm}$ do phần lồng vào nhau. Hỏi 2 mối lồng làm hao hụt tất cả bao nhiêu mi-li-mét?
\shortans{8}
\loigiai{
\textbf{Đáp số:} $8$\\
Tổng chiều dài 3 kẹp thẳng: $28 \times 3 = 84\text{ mm}$. Hao hụt: $84 - 76 = 8\text{ mm}$.
}
\end{ex}
